\secciongrande{Capacidades y limitaciones de ejecución ¿Por qué JavaScript en el cliente}

\resumen{
  JavaScript se utiliza en el lado cliente para mejorar la usabilidad, permitiendo interactividad sin acceder al sistema de archivos del servidor. Su ejecución está restringida por políticas de seguridad como el sandbox y la política del mismo origen, protegiendo al usuario, pero requiere validación de datos por parte del desarrollador.
}

\definicion{Lado cliente}{Entorno de ejecución del código (JavaScript) dentro del navegador del usuario.}
\definicion{Lado servidor}{Entorno donde se gestiona la lógica principal y se accede al sistema de archivos (ej. PHP, Python, Node.js).}
\subtituloA{Razones para usar JavaScript en el cliente}

\lista{
  **Navegabilidad sin JS:** El sitio debe ser funcional sin JavaScript; este último mejora la experiencia (rápida, moderna) pero no es esencial para la base. |
  **Actualización dinámica:** Examina datos recibidos del servidor y actualiza secciones de la página sin recarga completa.
}
